\section{Discussion}
\label{sec:discussion}

\subsection{Interpreting the 53\% Scale Confound}

The primary finding --- that MMLU scale control reduces the TruthfulQA $\times$ BBQ correlation by 53\% --- does not mean that factuality and bias avoidance are unrelated; the residual partial correlation ($0.343$, $p = 1.40 \times 10^{-9}$) is positive and highly significant. It means that the raw leaderboard correlation dramatically overstates the alignment-specific relationship.

What drives this? Models scoring high on MMLU are typically larger, more capable models trained on more data. Such models tend to answer more factual questions correctly (boosting TruthfulQA) and make fewer reasoning errors in bias-sensitive contexts (boosting BBQ), not necessarily because their alignment training is coherent, but because capability scales with both. This is the classic confounding structure: a common cause (scale) induces positive correlation between two effects, even when the direct relationship is weaker.

The implication for leaderboard interpretation is direct: a model family that improves TruthfulQA and BBQ together may simply be scaling up. Attributing such correlated improvement to aligned training objectives requires, at minimum, controlling for MMLU or an equivalent capability proxy.

\subsection{The Residual Coupling: What Does 0.343 Mean?}

After MMLU control, $\rho_{\text{partial}} = 0.343$ [0.180, 0.492] remains strongly significant. Three competing explanations deserve consideration:

\textbf{(1) Genuine RLHF co-training effect.} If RLHF and safety-focused fine-tuning jointly improve factuality and bias avoidance, cross-model variation in training intensity would induce residual coupling. However, our Tier~3 null result (RLHF sign test $p = 0.686$ on BBQ proxy) provides no cross-model evidence for this mechanism.

\textbf{(2) ARC Proxy artifact.} ARC Challenge measures logical reasoning; TruthfulQA also has a reasoning component. The partial correlation 0.343 may largely reflect shared reasoning demands, inflating the estimate beyond what genuine factuality $\times$ social-bias coupling would show. This is our preferred explanation.

\textbf{(3) Residual family-level clustering.} Within-family variation in training recipes may drive residual $\rho$ despite family-weighted correction.

Disambiguation requires genuine HELM Lite BBQ per-model accuracy. If $\rho_{\text{partial}}$ under real BBQ is substantially lower than 0.343, explanation (2) dominates.

\subsection{Limitations}

\textbf{L1: BBQ Proxy (ARC Challenge).} The most important limitation is the substitution of ARC Challenge accuracy for genuine HELM Lite BBQ per-model scores. The primary claim --- that MMLU confounds alignment benchmark correlations (Fisher $z$ $p = 3.22 \times 10^{-12}$) --- holds regardless of proxy choice; secondary claims about ``factuality-bias coupling'' are qualified by this proxy.

\textbf{L2: Scenario Classification Underpowered ($N{=}296$).} $\rho_{\text{partial}} = 0.343$ with CI [0.180, 0.492] spans the 0.40 structural threshold. Resolving the scenario requires $N \geq 400$--$600$ or genuine BBQ data. AMBIGUOUS is pre-registered and scientifically valid.

\textbf{L3: Safety Dimension Excluded (HarmBench $N{=}0$).} Three-benchmark partial Spearman matrix could not be computed. Requires alternative safety benchmark with better LLM LB v1 coverage.

\textbf{L4: RLHF Sign Test on Proxy.} Null result ($p = 0.686$) may reflect proxy inadequacy rather than genuine absence of RLHF bias effects.

\textbf{L5: Observational, Cross-Sectional, Open-Weight Only.} Results characterize $N{=}296$ open-weight LLMs from 2022--2024; do not generalize to proprietary models or post-2024 RLHF training.

\subsection{Implications for Alignment Evaluation Practice}

The Fisher $z$ difference test between raw and partial Spearman correlation should be a standard diagnostic step in alignment benchmark analysis whenever a capability proxy co-varies with the benchmarks of interest. More broadly, positive correlations between alignment benchmarks should not be interpreted as evidence of aligned constructs without first checking whether a common scale factor explains the co-movement.
